\subsection{自由阿贝尔群的群代数}

设 P 是一个有限秩自由 Z-模。记 A{[}P{]} 为 A 上加法群 P 的群代数（《代数》第三章第 2 节第 6 号）。对于任意 \(p \in P\)，记 \(e^{p}\) 为 A{[}P{]} 中相应的元素。于是 \((e^{p})_{p \in P}\) 是 A-模 A{[}P{]} 的一组基，并且对于任意 \(p, p' \in P\)，有

\[
e ^ {p} e ^ {p ^ {\prime}} = e ^ {p + p ^ {\prime}}, \quad (e ^ {p}) ^ {- 1} = e ^ {- p}, \quad e ^ {0} = 1.
\]

引理 1. 假设 A 是唯一分解环（《交换代数》第七章第 3 节第 1 号定义 1）。

\begin{enumerate}
\def\labelenumi{(\roman{enumi})}
\item
  环 A{[}P{]} 是唯一分解环。
\item
  若 \(u, v\) 是 \(\mathbf{P}\) 中不成比例的元素，则元素 \(1 - e^u, 1 - e^v\) 在 \(\mathrm{A}[\mathbf{P}]\) 中互素。
\end{enumerate}

令 \((p_{1}, p_{2}, \ldots, p_{l})\) 为 P 的一组基，并令 \(X_{1}, X_{2}, \ldots, X_{l}\) 为不定元。从 \(A[X_{1}, \ldots, X_{l}, X_{1}^{-1}, \ldots, X_{l}^{-1}]\) 到 A{[}P{]} 的 A-线性映射，将 \(X_{1}^{n_{1}} X_{2}^{n_{2}} \ldots X_{l}^{n_{l}}\)（其中 \(n_{1}, n_{2}, \ldots, n_{l} \in Z\)）映到 \(e^{n_{1} p_{1} + \cdots + n_{l} p_{l}}\)，它是环同构。现在 \(A[X_{1}, \ldots, X_{l}]\) 是唯一分解环（《交换代数》第七章第 3 节第 5 号），而 \(A[X_{1}, \ldots, X_{l}, X_{1}^{-1}, \ldots, X_{l}^{-1}]\) 是 \(A[X_{1}, \ldots, X_{l}]\) 的分式环，因而也是唯一分解环。

令 \(P'\)（分别为 \(P''\)）为 \(P\) 中满足某个倍数属于 \(Zu + Zv\)（分别属于 \(Zu\)）的元素所成的集合。则群 \(P/P'\) 与 \(P'/P''\) 无挠，因而存在 \(P''\) 在 \(P'\) 中的一个补，以及存在 \(P'\) 在 \(P\) 中的一个补。因此，存在一组基 \((z_1, z_2, \ldots, z_l)\)（它是 \(Z\)-模 \(P\) 的一组基）以及有理整数 \(j, m, n\)，使得 \(u = jz_1, v = mz_1 + nz_2, j > 0, n > 0\)。置 \(X_i = e^{z_i}\)（其中 \(1 \leqslant i \leqslant l\)），于是有 \(1 - e^u = 1 - X_1^j, 1 - e^v = 1 - X_1^m X_2^n\)。令 \(K\) 为 \(A\) 的分式域的代数闭包，从而可将 \(A[P]\) 认同为环 \(B = K[X_1, \ldots, X_l, X_1^{-1}, \ldots, X_l^{-1}]\) 的一个子环。对于任意 \(j\) 次单位根 \(z, 1 - zX_1\)，它在

\[
\mathrm{K} \left[ \mathrm{X} _ {1}, \dots , \mathrm{X} _ {l} \right];
\]

此外，由 \(1 - z\mathrm{X}_1\) 生成的理想不包含 \(\mathbf{X}_i\) 的任何单项式。由此得出，理想 \((1 - z\mathrm{X}_1)\mathrm{B}\) 是 \(\mathbf{B}\) 的高度为 1 的素理想（《交换代数》第七章第 1 节第 6 号），从而 \(1 - zX_{1}\) 是 B 中的不可约元。因此，\(1 - X_{1}^{j}\) 在 B 中的不可约因子具有 \(1 - zX_{1}\) 的形式。现在这些因子中没有一个整除 \(1 - X_{1}^{m}X_{2}^{n}\)（因为从 B 到 B 的同态 \(f\) 满足 \(f(X_{1}) = z^{-1}\)、\(f(X_{i}) = X_{i}\)（对 \(i \geqslant 2\)），且

\[
f (1 - z X _ {1}) = 0 \quad \text { and } \quad f (1 - X _ {1} ^ {m} X _ {2} ^ {n}) = 1 - z ^ {- m} X _ {2} ^ {n} \neq 0).
\]

因此，\(1 - X_1^j\) 与 \(1 - X_1^m X_2^n\) 在 B 中互素。所以，\(1 - X_1^j\) 与 \(1 - X_1^m X_2^n\) 在 A{[}P{]} 中的任一公因子在 B 中可逆，因而除以形如 \(X_1^{k_1}X_2^{k_2}\ldots X_l^{k_l}\) 的元素外，等于 A 中的某个元素 \(a\)；换言之，\(a\) 在 A 中整除 1，故在 A 中可逆。因此最后 \(1 - X_1^j\) 与 \(1 - X_1^m X_2^n\) 在 A{[}P{]} 中互素。

